\section{引言：规则从哪里来}\label{sec:intro}

\subsection{问题}

上一份综述得出的核心结论是：趋势跟踪与均值回归共享同一个信号，唯一区别是符号；两者的分野不在信号里，而在信号之后的每一步。这个结论对“设计规则”这件事提出了一个苛刻的要求——如果两策略的信号相同，那么规则设计的重心就\emph{不应该}放在“找一个更好的信号”上，而应该放在验证、缩放、对冲、交易与融资这五个环节上。

本文据此设计三套规则。选择这三套而不是别的，依据是论文实证中“哪几个口径曾经给出过正收益”：

\begin{itemize}[leftmargin=2.2em]
  \item \textbf{期货时序趋势}：论文表 5 显示，在中国期货时序口径下，趋势策略在 $L=5$ 到 $L=250$ 的全部六个窗口上净夏普为正（$0.09\sim0.32$），而时序反转全为负（$-0.17\sim-0.60$）。这是论文中趋势一侧唯一全窗口成立的口径。
  \item \textbf{A 股截面反转}：论文表 6 显示，A 股个股截面反转在 $L=20$ 到 $L=250$ 上净夏普正（$0.56\sim0.99$），其中 $L=60$ 净夏普 $0.99$、年化 $8.98\%$、最大回撤 $-18.64\%$、月胜率 $60\%$。这是论文全部实验中夏普最高的一组。
  \item \textbf{期货截面统计套利}：论文表 5 显示，期货截面反转在 $L=60$ 与 $L=120$ 上净夏普 $0.72$ 与 $0.56$，且 $L=60$ 的 Newey--West $t$ 统计量为 $2.59$——这是论文中少数通过常规显著性门槛的结果。
\end{itemize}

其余口径（A 股宽基指数时序）论文已明确标注证据力弱，本文不为它设计规则。

\subsection{三条设计铁律}

\subsubsection*{铁律一：每条规则必须能追溯到论文的具体证据}

交易规则最常见的失效方式不是逻辑错误，而是\emph{无据可依}——某个参数取 $60$ 日、某个阈值取 $2$ 倍标准差，说不出为什么。本文要求每套规则中的每个关键决策都指回论文中的具体表格或结论。

这条铁律带来一个副作用，必须提前说明：\textbf{论文没有检验过的东西，本文不得声称为有效}。例如多窗口信号集成、ATR 止损的具体倍数、参与率约束的冲击成本模型，这些在论文中或未检验、或仅作为执行章的定性讨论出现。本文会给出这些做法的工程建议，但会在证据映射表中把它们标为\emph{工程判断}而非\emph{实证支持}，与有论文支撑的条目明确分开。

\subsubsection*{铁律二：成本先于信号}

论文实证中最刺眼的一条证据是 A 股截面反转在 $L=5$ 上的表现：毛利夏普 $+1.08$，净夏普 $-0.35$。信号本身相当强（毛夏普超过 1），但 $79.8$ 倍的年化换手带来 $12.31\%$ 的年化成本拖累，把信号完全吞掉。

这条证据确立了本文的成本门槛：

\begin{quote}
\itshape
任何参数组合，只有当预期毛收益超过其往返成本的 3 倍时，才允许进入考察范围。
达不到这条线的组合，无论其信号统计量多显著，一律废弃。
\end{quote}

门槛设为 3 倍而非更高，是因为论文中表现最好的一档（期货截面反转 $L=60$）的成本占毛利比约为 $6.3\%$，即毛利约为成本的 16 倍；而 $L=60$ 的 A 股反转成本占毛利 $31.87\%$，约 3.1 倍——恰好在门槛线上。换言之，3 倍这条线是“论文实证中最勉强可用的组合”所对应的位置。它不是一个理论最优值，而是一个从实证中反推出来的下限。

\subsubsection*{铁律三：失效防线内置}

论文的对偶六指出两策略的失效模式不同：趋势策略是“震荡市慢性出血叠加崩溃时跳崖”，反转策略是“不收敛加资金断裂、跳变式亏损”。防线因此不能通用。

\begin{itemize}[leftmargin=2.2em]
  \item 趋势规则内置两道防线：\textbf{恐慌状态过滤}（对移动量崩溃）与\textbf{拥挤监控}（对信号同质化）。
  \item 反转规则内置两道防线：\textbf{偏离止损}（对不收敛）与\textbf{资金冗余}（对被迫平仓）。
\end{itemize}

需要特别强调一条反直觉的纪律：\textbf{反转策略不能用“越跌越买”的逻辑加仓}。论文在处理 LTCM 与 2007 年量化危机时指出，反转策略的亏损本身不可怕，可怕的是亏损触发了资金约束——业绩与资金流的正反馈会让套利者在错价最深时被迫减仓\cite{shleifervishny1997,khandanilo2011}。因此反转规则的加仓必须绑定在“价差已开始收敛”的证据上，而不是绑在“偏离更大了”上。

\subsection{三套规则的共同骨架}

本文对每套规则采用同一套叙述结构，便于横向比较与逐项复核：

\begin{enumerate}[leftmargin=2.2em]
  \item \textbf{适用性依据}：论文中支持该规则的具体证据；
  \item \textbf{信号}：计算什么、用什么窗口、如何处理噪声；
  \item \textbf{组合构建}：选多少标的、如何加权、如何中性化；
  \item \textbf{仓位与缩放}：波动率目标化是否适用（三套规则答案不同）；
  \item \textbf{进出场纪律}：调仓频率、免交易带、入场门槛；
  \item \textbf{止损与失效防线}：针对该策略失效模式的防线；
  \item \textbf{成本门槛与容量}：成本核算、容量约束来源；
  \item \textbf{参数边界表}：每个参数的取值、论文依据、越界后果；
  \item \textbf{证据映射表}：决策 $\to$ 论文结论 $\to$ 证据来源 $\to$ 置信度。
\end{enumerate}

第~\ref{sec:shared}~节先给出三套规则共享的基础设施（数据口径、成本模型、可交易性过滤、评估指标），此后三节各展开一套规则。

\subsection{边界声明}

本文是规则设计文档，不是策略推介。三条限定必须写在前面。

\textbf{其一，绝对收益不可预期}。论文的自建回测是变量受控的相对比较，未做多重检验校正（$6$ 个窗口 $\times$ 两个方向 $\times$ 四组口径 $=48$ 个组合），且 A 股侧做空不可实现、除权修复存在残余污染。本文沿用的所有绩效数字（如 $L=60$ 的 $8.98\%$ 年化）都应读作“在该口径下的相对位置”，不是可实现的收益预期。

\textbf{其二，A 股反转的多空版本不可直接执行}。论文披露 A 股融券受限、成本高、可借券池小，截面多空组合在多数时段不可实现。本文因此为 A 股反转给出两个版本：可执行的“纯多头加指数对冲”版本，与研究用的“多空”版本，后者明确标注不可实盘使用。

\textbf{其三，制度环境不可外推}。论文样本期为 2014--2026 年，其间中国期货与 A 股经历了股指期货交易限制、融券政策调整、注册制改革等制度变化。规则的适用性以该期间为限。
